% !TEX program = lualatex
% programmes-cpge.tex — documentation du paquet programmes-cpge
% Se compile sous LuaLaTeX comme sous pdfLaTeX, avec --shell-escape pour l'index
% (ou en lançant makeindex à la main).
\documentclass[a4paper,11pt]{article}

\usepackage{iftex}
\ifPDFTeX
  \usepackage[T1]{fontenc}
  \usepackage{lmodern}
\else
  \usepackage{fontspec}
\fi
\usepackage[french]{babel}
\usepackage[margin=2.4cm,marginparwidth=2cm]{geometry}
\usepackage{microtype}
\usepackage{imakeidx}
\makeindex[columns=2,title=Index des commandes]
\usepackage{tabularx,multicol,xcolor,amssymb}
\usepackage[hidelinks]{hyperref}

\usepackage{programmes-cpge}

% Les exemples sont en UTF-8 : listingsutf8 les recode en Latin-1 sous pdfLaTeX ;
% LuaLaTeX lit l'UTF-8 sans aide.
\ifPDFTeX
  \tcbuselibrary{documentation,listingsutf8}
  \newcommand\encodagelistings{inputencoding=utf8/latin1,}
\else
  \tcbuselibrary{documentation,listings}
  \newcommand\encodagelistings{}
\fi
\tcbset{
  doclang/.cd,
  color=couleur, colors=Couleurs, counter=compteur, counters=Compteurs,
  environment content=contenu de l'environnement, environment=environnement,
  environments=Environnements, commands=Commandes, path=chemin, paths=Chemins,
  index=Index, key=clé, keys=Clés, length=longueur, lengths=Longueurs,
  new=Nouveau, pageshort={p.}, updated=Mis à jour, value=valeur, values=Valeurs,
}
% Les options des exemples, communes aux deux moteurs ; l'encodage s'y ajoute.
\newcommand\optionsexemples{%
    style=tcbdocumentation, basicstyle=\ttfamily\small,
    extendedchars=true,
    moretexcs={cpgechamp,cpgetitre,cpgenumero,cpgeintitule,cpgegenre,cpgeparent,
      cpgefiliereDe,cpgeancetre,cpgenombre,cpgenombreenfants,cpgenombrelignes,cpgesiexiste,cpgesivide,
      cpgesiligne,cpgeafficher,cpgechapeau,cpgecapacites,cpgecite,cpgecettepartie,
      cpgepourchaque,cpgepourchaqueligne,cpgerangees,cpgetableau,cpgeboite,cpgecharger,
      cpgefiliere,cpgerubrique,cpgeligne,cpgeitem,cpgenotion,cpgevec,cpgegrad,cpgerot,
      cpgediv,cpgedd,cpgeii,cpgejj,cpgeRe,cpgerenvoi,cpgeformatcitation,
      cpgeformatexperimentale,cpgeformatnumerique,cpgeformattitretableau,
      cpgeformatnumeroligne,cpgeformattitreboite,tcbset},}
\edef\x{\noexpand\tcbset{documentation listing options={%
  \unexpanded\expandafter{\optionsexemples},\encodagelistings}}}
\x
\tcbset{
  documentation minted options={},
  color command=red!60!black, color key=blue!60!black, color value=teal!70!black,%
}

\newcommand\pkg[1]{\textsf{#1}}
\newcommand\cle[1]{\texttt{#1}}
\newcommand\champ[1]{\texttt{#1}}

\title{\pkg{programmes-cpge}\\[0.4em]
  \large Les programmes officiels de physique-chimie des CPGE,\\
  à citer, parcourir et composer en \LaTeX}
\author{Antonin Marchand}
\date{Version 0.2 — 21 octobre 2026}

\begin{document}
\maketitle

\begin{abstract}
Le paquet \pkg{programmes-cpge} met à disposition, sous forme de données \LaTeX,
le texte des programmes officiels de physique-chimie de cinq classes
préparatoires scientifiques — MPSI, PCSI, MP, PC et PSI — ligne par ligne : les
notions et contenus, les capacités exigibles, les capacités expérimentales et
numériques, les chapeaux de chaque partie. Chaque élément porte une clé stable
(\cle{MP/4/1/1} est la première ligne de l'électrostatique de MP). Le paquet
fournit de quoi lire un champ, parcourir une partie, et composer des tableaux
et des boîtes prêts à l'emploi, ou des mises en page sur mesure. Il fonctionne
sous pdf\LaTeX{} comme sous Lua\LaTeX.
\end{abstract}

\tableofcontents

% =========================================================================
\section{Ce que contient le paquet}

Chaque filière est un fichier de données, \texttt{programmes-cpge-<filière>.tex},
que le paquet lit au chargement.

\begin{center}
\begin{tabularx}{\linewidth}{@{}llXlr@{}}
  \toprule
  Clé & Discipline & Voie & Texte & Lignes \\
  \midrule
  \cle{MPSI} & \cpgechamp{MPSI}{discipline} & \cpgechamp{MPSI}{voie} & \cpgechamp{MPSI}{annee} & \cpgenombrelignes{MPSI} \\
  \cle{PCSI} & \cpgechamp{PCSI}{discipline} & \cpgechamp{PCSI}{voie} & \cpgechamp{PCSI}{annee} & \cpgenombrelignes{PCSI} \\
  \cle{MP}   & \cpgechamp{MP}{discipline}   & \cpgechamp{MP}{voie}   & \cpgechamp{MP}{annee}   & \cpgenombrelignes{MP} \\
  \cle{PC}   & \cpgechamp{PC}{discipline}   & \cpgechamp{PC}{voie}   & \cpgechamp{PC}{annee}   & \cpgenombrelignes{PC} \\
  \cle{PSI}  & \cpgechamp{PSI}{discipline}  & \cpgechamp{PSI}{voie}  & \cpgechamp{PSI}{annee}  & \cpgenombrelignes{PSI} \\
  \bottomrule
\end{tabularx}
\end{center}

La PCSI est la première année de la PC comme de la PSI. Pour la PCSI et la PC,
seul le programme de physique est fourni : leur programme de chimie est un texte
à part. L'année est celle du texte officiel, et non celle de sa mise en
application.

Chaque programme comprend, comme le texte du ministère, un préambule
(\cle{P}), une formation expérimentale (\cle{E}), les contenus thématiques
(\cle{1}, \cle{2}…) et trois annexes : matériel (\cle{A1}), outils mathématiques
(\cle{A2}) et outils numériques (\cle{A3}).

% =========================================================================
\section{Charger le paquet}

\begin{docCommand}{usepackage}{\oarg{filières}\texttt{\{programmes-cpge\}}}
  Charge les filières nommées, séparées par des virgules ; sans option, les
  cinq. L'option \cle{toutes} les charge explicitement.
\end{docCommand}

\begin{dispListing}
\usepackage[MP,PSI]{programmes-cpge}   % deux filières seulement
\usepackage{programmes-cpge}           % les cinq
\end{dispListing}

\begin{docCommand}{cpgecharger}{\marg{filières}}
  Charge des filières supplémentaires, n'importe où. Une filière déjà chargée ne
  l'est pas deux fois.
\end{docCommand}

Le paquet charge \pkg{amsmath}, \pkg{amssymb} (sauf si \pkg{unicode-math} est
déjà chargé), \pkg{array}, \pkg{booktabs}, \pkg{longtable} et \pkg{tcolorbox}.
Un préambule complet, pour chacun des deux moteurs :

\begin{dispListing}
% pdfLaTeX
\documentclass{article}
\usepackage[T1]{fontenc}
\usepackage{lmodern}
\usepackage[french]{babel}
\usepackage{programmes-cpge}

% LuaLaTeX
\documentclass{article}
\usepackage{fontspec}
\usepackage[french]{babel}
\usepackage{programmes-cpge}
\end{dispListing}

% =========================================================================
\section{Les clés}

Une clé est un chemin : la filière, puis les numéros du texte officiel, un cran
par niveau.

\begin{center}
\begin{tabular}{@{}ll@{}}
  \toprule
  Clé & Ce qu'elle désigne \\
  \midrule
  \cle{MP} & le programme de MP \\
  \cle{MP/4} & le thème 4 : \cpgetitre{MP/4} \\
  \cle{MP/4/1} & la partie 4.1 : \cpgetitre{MP/4/1} \\
  \cle{MP/4/1/1} & la première ligne de cette partie \\
  \cle{MPSI/1/7/2} & une sous-partie : \cpgetitre{MPSI/1/7/2} \\
  \cle{MP/E/1} & la formation expérimentale : \cpgetitre{MP/E/1} \\
  \cle{MP/P/3/APP} & une compétence du préambule : \cpgetitre{MP/P/3/APP} \\
  \cle{MPA/A2/1} & une annexe : \cpgetitre{MPA/A2/1} \\
  \bottomrule
\end{tabular}
\end{center}

Les annexes ont leur racine propre, \cle{MPA}, \cle{PCA}…, parce que le texte
officiel les numérote à part : \cle{MPA/A2/1/1} ne se confond pas avec
\cle{MP/2/1/1}. Elles restent des enfants de la filière quand on la parcourt.

Une clé désigne l'un de trois \emph{genres} d'éléments, chacun avec ses
\emph{champs} :

\begin{center}
\begin{tabularx}{\linewidth}{@{}lX@{}}
  \toprule
  Genre & Champs \\
  \midrule
  \cle{filiere} & \champ{discipline}, \champ{classe}, \champ{voie}, \champ{annee} \\
  \cle{rubrique} & \champ{parent}, \champ{numero}, \champ{niveau} (1 pour un thème,
    2 pour une partie, 3 pour une sous-partie), \champ{titre}, \champ{chapeau} \\
  \cle{ligne} & \champ{parent}, \champ{numero}, \champ{contenu}, \champ{capacites},
    \champ{experimentales}, \champ{numeriques} \\
  \bottomrule
\end{tabularx}
\end{center}

Tous portent aussi \champ{cle}, \champ{genre} et \champ{filiere}. Les champs
\champ{chapeau}, \champ{contenu} et les trois familles de capacités sont des
\emph{listes d'items} : un item par paragraphe du texte officiel, une capacité
exigible par item.

% =========================================================================
\section{Lire}

Les commandes de lecture sont \emph{développables} : on peut s'en servir dans un
titre de section, dans un \cs{edef}, dans un test.

\begin{docCommand}{cpgechamp}{\marg{clé}\marg{champ}}
  Le champ, tel quel. Une liste d'items se compose en paragraphes.
\end{docCommand}

\begin{docCommand}{cpgetitre}{\marg{clé}}\end{docCommand}
\begin{docCommand}{cpgenumero}{\marg{clé}}\end{docCommand}
\begin{docCommand}{cpgeintitule}{\marg{clé}}
  Le titre, le numéro, ou les deux à la suite.
\end{docCommand}

\begin{docCommand}{cpgegenre}{\marg{clé}}\end{docCommand}
\begin{docCommand}{cpgeparent}{\marg{clé}}\end{docCommand}
\begin{docCommand}{cpgefiliereDe}{\marg{clé}}
  Le genre (\cle{filiere}, \cle{rubrique} ou \cle{ligne}), la clé du parent, et
  la filière.
\end{docCommand}

\begin{docCommand}{cpgeancetre}{\marg{clé}\marg{niveau}}
  La clé de l'ancêtre au niveau voulu : le thème d'une ligne au niveau 1, sa
  partie au niveau 2. Rien si la clé n'a pas d'ancêtre à ce niveau.
\end{docCommand}

\begin{docCommand}{cpgenombre}{\marg{clé}\marg{champ}}\end{docCommand}
\begin{docCommand}{cpgenombreenfants}{\marg{clé}}
  Le nombre d'items d'un champ, et le nombre d'enfants d'un élément.
\end{docCommand}
\begin{docCommand}{cpgenombrelignes}{\marg{clés}}
  Le nombre de lignes sous les clés données, à toute profondeur. Elle compte en
  parcourant : elle n'est pas développable.
\end{docCommand}

\begin{dispExample}
\cpgeintitule{PC/5/2}, dans le thème
« \cpgetitre{\cpgeancetre{PC/5/2/1/4}{1}} » ;
la ligne \cpgenumero{PC/5/2/1/4} porte
\cpgenombre{PC/5/2/1/4}{capacites} capacités ;
le programme de PSI date de \cpgechamp{PSI}{annee}.
\end{dispExample}

\begin{docCommand}{cpgesiexiste}{\marg{clé}\marg{vrai}\marg{faux}}\end{docCommand}
\begin{docCommand}{cpgesivide}{\marg{clé}\marg{champ}\marg{vrai}\marg{faux}}\end{docCommand}
\begin{docCommand}{cpgesiligne}{\marg{clé}\marg{vrai}\marg{faux}}
  Les tests : la clé existe-t-elle, le champ est-il vide, la clé est-elle une
  ligne.
\end{docCommand}

\begin{dispExample}
\cpgesiexiste{MP/9/9}{existe}{n'existe pas} ;
la ligne MP/2/2/1 \cpgesivide{MP/2/2/1}{capacites}%
  {n'a que des capacités expérimentales}{a des capacités exigibles}.
\end{dispExample}

% =========================================================================
\section{Composer un champ}

\begin{docCommand}{cpgeafficher}{\oarg{options}\marg{clé}\marg{champ}}
  Compose un champ selon une forme :
  \begin{docKey}{forme}{=\textcolor{blue!60!black}{\texttt{paragraphes}}\,|\,\texttt{liste}\,|\,\texttt{enumeration}\,|\,\texttt{ligne}}{par défaut \texttt{paragraphes}}
    Un item par paragraphe, en liste à puces, en liste numérotée, ou à la suite.
  \end{docKey}
  \begin{docKey}{separateur}{=\meta{texte}}{par défaut une espace}
    Ce qui sépare deux items de la forme \texttt{ligne}.
  \end{docKey}
  \begin{docKey}{format}{=\meta{commande}}{par défaut aucune}
    Une commande à un argument appliquée à chaque item.
  \end{docKey}
\end{docCommand}

\begin{dispExample}
\cpgeafficher[forme=liste]{MP/4/1/4}{capacites}
\cpgeafficher[forme=ligne, separateur={ ; },
  format=\emph]{MP/4/1/4}{contenu}
\end{dispExample}

\begin{docCommand}{cpgechapeau}{\oarg{options}\marg{clé}}
  Le chapeau d'une rubrique. Il cite souvent sa partie : \cs{cpgecettepartie}
  y désigne alors la rubrique composée.
\end{docCommand}

\begin{dispExample}
\cpgechapeau{PSI/4/6}
\end{dispExample}

\begin{docCommand}{cpgecapacites}{\oarg{options}\marg{clé}}
  Les capacités d'une ligne, les trois familles à la suite : exigibles,
  expérimentales (en gras, comme dans le texte officiel), numériques (précédées
  de « Capacité numérique : »).
  \begin{docKey}{forme}{=\texttt{paragraphes}\,|\,\texttt{liste}}{par défaut \texttt{paragraphes}}\end{docKey}
  \begin{docKey}{exigibles}{=\texttt{true}\,|\,\texttt{false}}{par défaut \texttt{true}}\end{docKey}
  \begin{docKey}{experimentales}{=\texttt{true}\,|\,\texttt{false}}{par défaut \texttt{true}}\end{docKey}
  \begin{docKey}{numeriques}{=\texttt{true}\,|\,\texttt{false}}{par défaut \texttt{true}}
    Les familles à composer.
  \end{docKey}
\end{docCommand}

\begin{dispExample}
\cpgecapacites[forme=liste]{MPSI/1/5/3}
\end{dispExample}

\begin{docCommand}{cpgecite}{\marg{clé}}\end{docCommand}
\begin{docCommand}{cpgecettepartie}{}
  Le numéro et le titre d'une rubrique, mis en forme par
  \cs{cpgeformatcitation}. \cs{cpgecettepartie} cite la rubrique en cours de
  composition, et vaut « cette partie » hors de tout contexte. Comme toute
  commande sans argument, elle mange l'espace qui la suit : on écrit
  \verb+\cpgecettepartie{} constitue+.
\end{docCommand}

\begin{dispExample}
Voir la partie \cpgecite{PCSI/1/5}.
\end{dispExample}

% =========================================================================
\section{Parcourir}

\begin{docCommand}{cpgepourchaque}{\marg{clé}\marg{code}}
  Exécute le code pour chaque enfant de la clé, dans l'ordre du texte ;
  \texttt{\#1} y vaut la clé de l'enfant.
\end{docCommand}

\begin{docCommand}{cpgepourchaqueligne}{\marg{clés}\marg{code}}
  Exécute le code pour chaque \emph{ligne} sous les clés données, à toute
  profondeur ; une clé de ligne se désigne elle-même. Les clés se séparent par
  des virgules.
\end{docCommand}

\begin{dispExample}
\cpgepourchaque{PC}{\cpgeintitule{#1}\par}
\end{dispExample}

\begin{dispExample}
Les lignes de \cpgecite{MP/4/4} :
\cpgepourchaqueligne{MP/4/4}{\cpgenumero{#1}~}.
\end{dispExample}

\begin{docCommand}{cpgerangees}{\marg{clés}\marg{modèle}}
  Les rangées d'un tableau fait main : le modèle, où \texttt{\#1} vaut la clé,
  est répété pour chaque ligne sous les clés données. Les rangées sont
  assemblées avant d'être rendues : les \texttt{\&} et \cs{\textbackslash} du
  modèle entrent dans l'alignement comme s'ils étaient écrits à la main, dans
  \texttt{tabular}, \texttt{tabularx} ou \texttt{longtable}.
\end{docCommand}

\begin{dispExample}
\begin{tabularx}{\linewidth}{@{}lX@{}}
  \toprule
  \cpgerangees{MP/4/3}{\cpgenumero{#1} &
    \cpgeafficher[forme=ligne]{#1}{contenu} \\}
  \bottomrule
\end{tabularx}
\end{dispExample}

% =========================================================================
\section{Mises en page prêtes à l'emploi}

\subsection{Le tableau du texte officiel}

\begin{docCommand}{cpgetableau}{\oarg{options}\marg{clés}}
  Le tableau à deux colonnes du texte officiel, « Notions et contenus » et
  « Capacités exigibles », pour les clés données : une ligne, une partie, un
  thème, plusieurs à la fois. Une rubrique y porte son titre sur toute la
  largeur. C'est un \texttt{longtable} : il se coupe d'une page à l'autre et
  répète son en-tête.
  \begin{docKey}{largeur}{=\meta{nombre}}{par défaut \texttt{0.38}}
    La part de la largeur de ligne donnée à la première colonne.
  \end{docKey}
  \begin{docKey}{entete}{=\texttt{true}\,|\,\texttt{false}}{par défaut \texttt{true}}\end{docKey}
  \begin{docKey}{titres}{=\texttt{true}\,|\,\texttt{false}}{par défaut \texttt{true}}
    Les titres des rubriques traversées.
  \end{docKey}
  \begin{docKey}{numeros}{=\texttt{true}\,|\,\texttt{false}}{par défaut \texttt{false}}
    La clé de chaque ligne, en petit, au-dessus de son contenu.
  \end{docKey}
  \begin{docKey}{exigibles}{=\texttt{true}\,|\,\texttt{false}}{par défaut \texttt{true}}\end{docKey}
  \begin{docKey}{experimentales}{=\texttt{true}\,|\,\texttt{false}}{par défaut \texttt{true}}\end{docKey}
  \begin{docKey}{numeriques}{=\texttt{true}\,|\,\texttt{false}}{par défaut \texttt{true}}\end{docKey}
  \begin{docKey}{long}{=\texttt{true}\,|\,\texttt{false}}{par défaut \texttt{true}}
    Un \texttt{tabular} au lieu d'un \texttt{longtable} : nécessaire dans une
    boîte, dans deux colonnes, partout où \texttt{longtable} ne va pas.
  \end{docKey}
  \begin{docKey}{titre contenus}{=\meta{texte}}{par défaut \texttt{Notions et contenus}}\end{docKey}
  \begin{docKey}{titre capacites}{=\meta{texte}}{par défaut \texttt{Capacités exigibles}}\end{docKey}
\end{docCommand}

Une partie entière, comme elle paraît dans le texte officiel :

\begin{dispListing}
\cpgetableau{MP/2/1}
\end{dispListing}

\cpgetableau{MP/2/1}

Deux lignes choisies, sans titre, avec leur clé :

\begin{dispExample}
\cpgetableau[long=false, titres=false, numeros, largeur=0.33]
  {PSI/4/6/1, PSI/4/6/5}
\end{dispExample}

\subsection{La boîte}

\begin{docCommand}{cpgeboite}{\oarg{options}\marg{clé}}
  Une ligne en boîte : son contenu au-dessus, ses capacités en liste
  au-dessous. Une rubrique en boîte : son chapeau, et ses lignes si on le
  demande.
  \begin{docKey}{lignes}{=\texttt{true}\,|\,\texttt{false}}{par défaut \texttt{false}}\end{docKey}
  \begin{docKey}{chapeau}{=\texttt{true}\,|\,\texttt{false}}{par défaut \texttt{true}}\end{docKey}
  \begin{docKey}{titre}{=\meta{texte}}{par défaut le numéro et le titre}\end{docKey}
  \begin{docKey}{tcolorbox}{=\meta{options}}{par défaut aucune}
    Des options \pkg{tcolorbox} pour cette boîte.
  \end{docKey}
  \begin{docKey}{exigibles}{=\texttt{true}\,|\,\texttt{false}}{par défaut \texttt{true}}\end{docKey}
  \begin{docKey}{experimentales}{=\texttt{true}\,|\,\texttt{false}}{par défaut \texttt{true}}\end{docKey}
  \begin{docKey}{numeriques}{=\texttt{true}\,|\,\texttt{false}}{par défaut \texttt{true}}
    Les familles de capacités à composer, comme pour \cs{cpgecapacites}.
  \end{docKey}
\end{docCommand}

\begin{dispExample}
\cpgeboite{MP/4/1/4}
\end{dispExample}

\begin{dispExample}
\cpgeboite[lignes, tcolorbox={colframe=teal!60!black,
  colbacktitle=teal!10}]{PC/1/3}
\end{dispExample}

Toutes les boîtes suivent le style \pkg{tcolorbox} \texttt{cpge boite}, qu'on
retouche une fois pour le document :

\begin{dispListing}
\tcbset{cpge boite/.append style={colframe=blue!50!black,
  colbacktitle=blue!8, fonttitle=\sffamily\bfseries}}
\end{dispListing}

% =========================================================================
\section{Galerie : des mises en page sur mesure}

Les commandes de lecture et de parcours suffisent à toutes les mises en page.
Voici quelques recettes, à recopier et à adapter.

\subsection{Une fiche d'auto-évaluation}

Chaque capacité d'une partie précédée d'une case, par l'option \texttt{format},
dans une boîte \pkg{tcolorbox} ordinaire :

\begin{dispExample}
\newcommand\caseacocher[1]{$\square$~#1}
\begin{tcolorbox}[title={Je sais faire -- \cpgeintitule{MP/4/4}}]
  \cpgepourchaqueligne{MP/4/4}{%
    \cpgeafficher[forme=liste, format=\caseacocher]{#1}{capacites}}
\end{tcolorbox}
\end{dispExample}

\subsection{Un tableau de progression}

Une colonne pour la date à laquelle chaque ligne a été traitée, à remplir à la
main :

\begin{dispExample}
\begin{tabularx}{\linewidth}{@{}lXp{2cm}@{}}
  \toprule
  Ligne & Notions et contenus & Traitée le \\
  \midrule
  \cpgerangees{PCSI/1/5}{\cpgenumero{#1} &
    \cpgeafficher[forme=ligne]{#1}{contenu} & \\ \midrule}
\end{tabularx}
\end{dispExample}

\subsection{Les capacités expérimentales d'un thème}

Un filtre : seules les lignes qui portent une capacité expérimentale, et
seulement celle-ci.

\begin{dispExample}
\cpgepourchaqueligne{PSI/1/4}{%
  \cpgesivide{#1}{experimentales}{}{%
    \cpgeboite[exigibles=false, numeriques=false,
      tcolorbox={colback=orange!4, colframe=orange!60!black}]{#1}}}
\end{dispExample}

\subsection{Le sommaire d'un programme}

\begin{dispExample}
\begin{multicols}{2}\small
\cpgepourchaque{MP}{%
  \textbf{\cpgeintitule{#1}}\par
  \cpgepourchaque{#1}{%
    \cpgesiligne{##1}%
      {\quad\textit{\cpgenumero{##1}
        \cpgeafficher[forme=ligne]{##1}{contenu}}\par}%
      {\quad\cpgeintitule{##1}\par}}}
\end{multicols}
\end{dispExample}

Dans un code imbriqué, la clé de l'enfant s'écrit \texttt{\#\#1}, comme dans
toute définition imbriquée. Une rubrique a pour enfants des sous-rubriques
ou des lignes : les annexes du matériel (A1) et des outils numériques (A3)
n'ont que des lignes, sans titre. Le test \cs{cpgesiligne} les distingue, et
l'exemple compose une ligne par son numéro et son contenu, en italique.

\subsection{Trois colonnes, une par famille de capacités}

\begin{dispExample}
\footnotesize
\begin{tabularx}{\linewidth}{@{}>{\raggedright}X>{\raggedright}X
    >{\raggedright\arraybackslash}X@{}}
  \toprule
  Contenus & Exigibles & Expérimentales \\
  \midrule
  \cpgerangees{MP/2/2}{%
    \cpgeafficher{#1}{contenu} & \cpgeafficher{#1}{capacites} &
    \cpgeafficher{#1}{experimentales} \\ \midrule}
\end{tabularx}
\end{dispExample}

% =========================================================================
\section{Les macros des données}

Le texte des programmes emploie quelques macros, que le paquet définit et que
l'on peut redéfinir par \cs{renewcommand} pour les accorder au reste d'un
document.

\begin{center}
\begin{tabularx}{\linewidth}{@{}lXl@{}}
  \toprule
  Macro & Rôle & Défaut \\
  \midrule
  \cs{cpgenotion}\marg{titre} & titre de notion, en tête d'un contenu & \cs{textbf} \\
  \cs{cpgecettepartie}, \cs{cpgecite} & citation d'une partie & voir plus haut \\
  \cs{cpgeformatcitation}\marg{n}\marg{t} & sa mise en forme & \textbf{n}~«~t~» \\
  \cs{cpgerenvoi}\marg{clé} & renvoi d'une capacité expérimentale à sa ligne & rien \\
  \cs{cpgevec}\marg{x} & vecteur & \cs{overrightarrow} \\
  \cs{cpgegrad}, \cs{cpgerot}, \cs{cpgediv} & opérateurs & $\cpgegrad$, $\cpgerot$, $\cpgediv$ \\
  \cs{cpgedd} & le d droit d'une différentielle & $\cpgedd x$ \\
  \cs{cpgeii}, \cs{cpgejj} & unités imaginaires & $\cpgeii$, $\cpgejj$ \\
  \cs{cpgeRe} & partie réelle & $\cpgeRe$ \\
  \bottomrule
\end{tabularx}
\end{center}

\begin{sloppypar}
Les mises en forme des mises en page prêtes à l'emploi se redéfinissent de
même : \cs{cpgeformatexperimentale}, \cs{cpgeformatnumerique},
\cs{cpgeformattitretableau}, \cs{cpgeformatnumeroligne},
\cs{cpgeformattitreboite}.
\end{sloppypar}

\begin{dispListing}
\usepackage{esvect}
\renewcommand\cpgevec{\vv}
\renewcommand\cpgenotion[1]{\textsc{#1}}
\renewcommand\cpgeformatnumerique[1]{\textsf{Python :} #1}
\end{dispListing}

% =========================================================================
\section{Ajouter ou corriger un programme}

Les fichiers \texttt{programmes-cpge-<filière>.tex} sont la source : c'est là
qu'on corrige une coquille, et c'est sur leur modèle qu'on ajoute une filière.
Le nom du fichier est en minuscules ; l'option du paquet est la clé de la
filière, en majuscules.

\begin{docCommand}{cpgefiliere}{\marg{clé}\marg{champs}}\end{docCommand}
\begin{docCommand}{cpgerubrique}{\marg{clé}\marg{champs}}\end{docCommand}
\begin{docCommand}{cpgeligne}{\marg{clé}\marg{champs}}
  Les trois déclarations d'un fichier de données. Les champs s'écrivent
  \texttt{nom = \{valeur\}}, séparés par des virgules ; un parent se déclare
  avant ses enfants, qui se rangent dans l'ordre du fichier.
\end{docCommand}

\begin{docCommand}{cpgeitem}{\marg{texte}}
  Un item d'un champ liste.
\end{docCommand}

\begin{dispListing}
\cpgefiliere{TSI}{
  discipline = {Physique-chimie}, classe = {TSI},
  voie = {Technologie et sciences industrielles}, annee = {2021},
}
\cpgerubrique{TSI/1}{
  parent = {TSI}, numero = {1}, niveau = {1}, titre = {Électronique},
  chapeau = {\cpgeitem{Premier paragraphe.}\cpgeitem{Second paragraphe.}},
}
\cpgeligne{TSI/1/1}{
  parent = {TSI/1}, numero = {1.1},
  contenu = {\cpgeitem{Signaux périodiques.}},
  capacites = {\cpgeitem{Définir la valeur moyenne d'un signal.}},
  experimentales = {}, numeriques = {},
}
\end{dispListing}

% =========================================================================
\section{pdf\LaTeX{} et Lua\LaTeX}

Le paquet est écrit en expl3 et ne dépend d'aucun moteur. Deux précautions :
\begin{itemize}
  \item sous pdf\LaTeX, les données sont en UTF-8 et emploient les guillemets
    français et l'apostrophe typographique : \verb|\usepackage[T1]{fontenc}|
    est nécessaire ;
  \item sous Lua\LaTeX, on ne charge pas \pkg{fontenc}, mais \pkg{fontspec}.
    Avec \pkg{unicode-math}, \pkg{amssymb} n'est pas chargé, ses symboles étant
    déjà définis.
\end{itemize}
Cette documentation se compile sous les deux moteurs.

% =========================================================================
\section{Les textes et leur licence}

Les textes des programmes sont ceux publiés par le ministère de l'Enseignement
supérieur et de la Recherche (\url{https://www.enseignementsup-recherche.gouv.fr}),
annexes des arrêtés de 2020 (première année) et de 2021 (seconde année). Ils
ont été relevés dans les documents officiels, et les formules mises en \LaTeX.
Une erreur de transcription se signale à l'auteur.

Ces textes sont des informations publiques. Les programmes de physique-chimie
ne portent aucune licence particulière --- contrairement à ceux de
mathématiques et d'informatique, publiés sous licence Creative Commons ---, si
bien que s'applique le régime par défaut : leur réutilisation est libre
(code des relations entre le public et l'administration, art.~L.~321-1), à
condition de ne pas les altérer, de ne pas en dénaturer le sens, et d'en citer
la source et la date (art.~L.~322-1). C'est ce que fait ce paquet : le texte est
transcrit tel quel, seule sa mise en forme est celle de \LaTeX, et chaque
filière porte l'année de son texte (\cs{cpgechamp}\marg{filière}\marg{annee}).

\subsection*{Les programmes composés}

Le dossier \texttt{doc} contient aussi un programme par filière, composé en
entier par le paquet : \texttt{programme-mpsi.pdf}, \texttt{programme-pcsi.pdf},
\texttt{programme-mp.pdf}, \texttt{programme-pc.pdf} et
\texttt{programme-psi.pdf}. Chaque ligne y porte sa clé. Ils se lisent à côté
du texte publié par le ministère : c'est la preuve de la concordance, et le
moyen le plus simple de trouver la clé d'une ligne. Leur source,
\texttt{programme.tex}, lit la filière dans le nom du job :
\begin{verbatim}
lualatex -jobname=programme-mp programme.tex
\end{verbatim}

\subsection*{Écarts au texte officiel}

Le texte est transcrit tel que le ministère l'a publié, coquilles comprises,
sauf deux coquilles évidentes du programme de MPSI, corrigées parce qu'elles ne
changent rien au sens. Chacune est signalée par un commentaire dans le fichier
de données.
\begin{itemize}
  \item Annexe 3 (outils numériques), chapeau : le texte officiel écrit « la
    formation des étudiants inclue » ; le paquet écrit « inclut ».
  \item Formation expérimentale, prévention du risque au laboratoire : le texte
    officiel numérote « 1. Prévention des risques au laboratoire », puis « 4.
    Prévention de l'impact environnemental » ; le paquet numérote la seconde
    « 2. », comme les programmes des autres filières.
\end{itemize}
La mise en forme, elle, est celle de \LaTeX{} : les formules sont composées,
« densité*vitesse » s'écrit $\text{densité}\times\text{vitesse}$, et une liste
du texte officiel reste une liste.

Le code du paquet est distribué sous la licence LPPL, version 1.3c ou
ultérieure.

% =========================================================================
\section{Historique}

\begin{description}
  \item[0.2 (2026-10-21)] Corrections.
    \begin{itemize}
      \item Le programme de MPSI compile de nouveau sous pdfLaTeX (une
        apostrophe dans une formule).
      \item La clé d'un élément se lit par \cs{cpgechamp}\marg{clé}\marg{cle},
        et l'option \texttt{numeros} de \cs{cpgetableau} l'affiche.
      \item L'espace après \cs{cpgecettepartie} n'est plus mangé dans les
        chapeaux.
      \item Les rubriques sont dans l'ordre du texte officiel, et le texte lui
        est conforme mot à mot, aux écarts près de la section « Les textes et
        leur licence ». Dans « Communiquer », en MP et MPSI, les clés
        \texttt{COM/3} et \texttt{COM/4} disparaissent : la capacité et sa
        sous-liste sont une seule ligne, comme dans l'officiel, et l'ancienne
        \texttt{COM/4} devient \texttt{COM/2}.
      \item Tirets et points de suspension en Unicode, au même rendu.
      \item Un programme composé par filière, dans \texttt{doc}.
    \end{itemize}
  \item[0.1 (2026-10-17)] Première version : MPSI, PCSI, MP, PC (physique) et
    PSI.
\end{description}

\printindex
\end{document}
